\newpage
\vspace{-1cm}
\chapter*{\zihao{-3}\heiti{ABSTRACT}}
\addcontentsline{toc}{chapter}{Abstract}
\vspace{-0.5cm}

Road networks form the core infrastructure of urban transportation systems, and their representation learning is critical for intelligent transportation applications such as route planning, travel time estimation, and traffic speed inference. Real-world road networks exhibit significant hierarchical structural properties: from individual road segments to structural localities and functional regions, forming a natural tree-like hierarchical organization with power-law degree distributions. However, existing Euclidean-space-based road network representation methods have fundamental limitations in encoding such hierarchical structures---the polynomial volume growth of Euclidean space is fundamentally mismatched with the exponential branching patterns of road network hierarchies, leading to severe embedding distortion and difficulty in precisely modeling hierarchical relationships.

Hyperbolic geometry, with its negative curvature and exponential volume growth, provides a natural geometric prior for embedding hierarchical data. This thesis adopts hyperbolic geometry as a unified tool and conducts progressive research on road network hierarchical structure modeling: from implicitly capturing hierarchical information through geometric properties of hyperbolic space, to explicitly constraining hierarchical relationships via entailment cone mechanisms, gradually deepening the geometric modeling capability for road network hierarchies.

This thesis first presents a systematic theoretical analysis of the compatibility between road network hierarchical properties and hyperbolic space. Through large-scale $\delta$-hyperbolicity analysis across 50 cities in 29 countries, the tree-like nature of road network topology is empirically verified. The theoretical advantages of hyperbolic space as a geometric framework for road network hierarchical modeling are demonstrated from three perspectives: exponential volume growth, hierarchical depth encoding, and parent-child relationship modeling. Based on this analysis, two paradigms for hierarchical modeling---implicit capture and explicit constraint---are proposed, providing a unified theoretical foundation for the subsequent two contributions.

The first contribution proposes the Hyperbolic Road Graph Neural Network (HRGNN), adopting the implicit hierarchical modeling paradigm. HRGNN employs the Poincar\'{e} ball model to capture the hierarchical structure of road networks, using a dual-pathway hyperbolic attention mechanism to simultaneously capture local neighborhood connectivity patterns and hierarchical dependencies revealed by importance-guided depth-first search (DFS) paths. An adaptive fusion strategy dynamically adjusts the contribution weights of the two pathways based on road network characteristics. Experiments on road type classification across multiple real-world city networks demonstrate that HRGNN improves micro-averaged F1 scores by 8.4\%--9.1\% over Euclidean baselines on intra-city classification, and improves macro-averaged F1 by 13.0\% over existing hyperbolic methods on cross-city transfer learning, validating the effectiveness of hyperbolic geometry for modeling road network hierarchies.

The second contribution proposes HHRoad, a hyperbolic hierarchical road network representation model with compositional entailment, adopting the explicit hierarchical modeling paradigm to further address the lack of geometric constraints on hierarchical relationships in implicit modeling. HHRoad operates in the Lorentz hyperbolic model, constructing a three-level hierarchy of road Segments, structural Localities, and functional Regions. Three core innovations enable explicit hierarchical constraints: (1) a hyperbolic embedding layer that maps road segment attributes onto the Lorentz hyperboloid, where distance from the origin naturally encodes hierarchical depth; (2) compositional entailment learning, which uses entailment cones to explicitly enforce the Region$\to$Locality$\to$Segment hierarchy; and (3) hyperbolic contrastive learning based on Lorentzian distance that captures both intra-level and cross-level semantic similarity. Experiments on five real-world city datasets within the VecCity benchmark framework show that HHRoad achieves the best overall average ranking (3.07) among eight methods, with ACC@3 improvements of up to 9.4\% over the strongest baselines on similarity trajectory search tasks.

The two contributions in this thesis progressively demonstrate that hyperbolic geometry provides an effective geometric framework for modeling road network hierarchical structures---from implicitly leveraging spatial properties to explicitly designing geometric constraints---naturally adapting to the tree-like topology of road networks, encoding hierarchical relationships with lower distortion, and yielding performance improvements across tasks and cities.

\vspace{0.5cm}
\hspace{-1cm}
{\bf{\sihao{Keywords:}}} \textit{Road Network Representation Learning; Hyperbolic Geometry; Graph Neural Network; Hierarchical Structure Modeling; Entailment Cones; Contrastive Learning}